\section{Language Models (Mostly) Know What They Know}
\paperhead{Language Models (Mostly) Know What They Know}{Saurav Kadavath, Tom Conerly, Amanda Askell, Tom Henighan, Dawn Drain et al.}{arXiv preprint (Anthropic), 2022}{https://arxiv.org/abs/2207.05221}{arXiv API id\_list=2207.05221, title match 1.00}

\subsection{What the paper does}

An honest system has to report its own competence, so the paper separates two questions: can a model judge whether a claim it just produced is correct, and can it predict in advance which questions it will get right at all. The starting point is plain calibration. When large models see multiple choice questions with visible letter-labelled options and are asked only for the letter, the probability mass they place on a letter tracks how often that letter is actually correct, and the agreement improves with model size and with few-shot prompting.

Calibration is then bootstrapped into self-knowledge via two distinct signals. P(True) is post-hoc: a candidate answer already exists and the model scores its validity. P(IK), the probability that ``I know'' the answer, is pre-hoc: it is read from the question alone, before any answer is proposed. The second is the harder object, because it is a claim about the model itself rather than about the world.

\subsection{How the method works}

Calibration is summarised by binning predictions into 10 equally populated bins and reporting expected and RMS calibration error. Format is load-bearing: the options must be visible and letter-labelled. Replacing one option with ``none of the above'' significantly degrades both accuracy and calibration, and 20-shot prompting does not repair it, so the authors reformulate every judgement as a True/False question of the shape ``Question: \ldots{} Proposed Answer: \ldots{} Is the proposed answer: (A) True (B) False'', and read the probability assigned to ``(A)''.

P(True) applies that template to the model's own generated samples on TriviaQA, Lambada, arithmetic, GSM8k and Codex HumanEval. A stronger variant first lists five of the model's own $T=1$ samples in the context as brainstormed ideas, then asks it to judge one designated possible answer. No finetuning is involved, only prompting; evaluation is 20-shot.

P(IK) instead adds a binary classification value head on top of the language model, read at the final token, with the whole model plus head finetuned under cross-entropy. The ground-truth label is the fraction of 30 unit-temperature samples of that question that are correct, approximated by repeated hard labels: a question answered correctly 20 times out of 30 contributes 20 ``IK'' copies and 10 ``IDK'' copies. Models of 800M, 3B, 12B and 52B parameters are studied.

\subsection{What it reports}

Large models are well calibrated in the letter-labelled format across BIG Bench, MMLU, TruthfulQA, LogiQA and QuALITY, with calibration improving as models scale and as prompting moves from 0-shot to 5-shot. Self-evaluation is far weaker zero-shot: P(True) sits close to 50\% for typical samples and is poorly calibrated, but improves few-shot, and showing five comparison samples improves Brier scores clearly on short-answer tasks and less on long-form ones. Filtering to samples with P(True) above 0.5 lifts task accuracy well above base accuracy, and the gap widens with model size.

P(IK) trained only on TriviaQA reaches AUROC 0.864 and Brier 0.151 in distribution, and transfers unevenly: 0.928 / 0.194 on mixed arithmetic, 0.687 / 0.164 on Python function synthesis, 0.624 / 0.200 on GSM8K, and 0.606 / 0.431 on LAMBADA. Training on all tasks except GSM8K improves every entry, moving LAMBADA to 0.853 / 0.108 and GSM8K to 0.752 / 0.121. P(IK) also responds to context it was never trained on: an obscure trivia question scores 18\% for the 52B model and 78\% once the relevant Wikipedia article is prepended, with larger gains for shorter articles, while correct hints raise P(IK) more than model-generated bad hints and hints copied from other questions push it down. RLHF policies appear badly miscalibrated until their distribution is tempered at $T=2.5$.

\subsection{How Mytheus uses this}
\begin{takeaway}
P(IK) is the closest published analogue of the \texttt{ContextGapAgent} trigger: not ``is this answer wrong'' but ``is something missing here'', which is exactly the condition that should fire an \texttt{agent-requests} message down the peer, related-meeting, RAG, silence chain. Two constraints block a literal copy. First, P(IK) is a finetuned value head, and Mytheus serves Anthropic models as a hosted API through LiteLLM, so there are no weights to attach a head to and no per-tenant head without breaking the isolation model. Second, the strongest P(True) configuration costs five extra unit-temperature samples inside a 20-shot context, which survives neither the p95 1.5 second peer round trip nor the Haiku triage cost cascade. What does transfer is the format discipline plus the source-material result. Build a single-sample prompted P(True) probe over the assembled RAG context and use the probability mass on the ``(A)'' token as the numeric confidence feeding the 0.7 surface gate and the 0.4 discard rule, replacing any free-text self-reported score; the fact that P(IK) rises sharply when the answer is genuinely present in context is direct evidence such a probe can distinguish a missing-context gap from a merely wrong answer. Reserve the five-sample variant for the off-critical-path \texttt{ProactiveAgent} conflict check on Opus. Measure rather than inherit: the golden-set harness already computes Brier, so add per-agent reliability curves over 10 equally populated bins, and cite the near-50\% zero-shot result as the reason the versioned prompt templates must ship few-shot exemplars. None of these numbers apply to free-form JSON extraction graded against role labels.
\end{takeaway}

\subsection{Limitations and what it leaves open}
\begin{itemize}[leftmargin=*, itemsep=1pt, topsep=2pt]
  \item Every evaluation is a static, single-turn question with one checkable answer, graded by string comparison or a test suite; calibration is never measured for structured multi-field output (a decision with rationale, stakeholders and confidence) nor for streaming, multi-speaker, incrementally revised input, so the 0.7 and 0.4 gates cannot inherit these curves.
  \item P(IK) requires model weights and an offline corpus of 30 unit-temperature samples per question. There is no hosted-API variant, no per-tenant head, and no discussion of what happens when the head is trained on one organisation's data and served to another.
  \item P(IK) generalises only partially and is badly calibrated out of distribution: Brier 0.431 on LAMBADA when trained on TriviaQA alone, versus 0.108 once LAMBADA is in the training mix, so a gap detector tuned on one meeting corpus has no established reliability on a new tenant.
  \item No latency, throughput or serving-cost figure appears anywhere, and the best-performing setting deliberately adds sampling and a 20-shot context, so the paper never asks whether self-evaluation is affordable inside a real-time budget.
  \item Cross-model experiments show only around a 6\% P(IK) difference between two models on the questions each answers correctly, and the calibration results hold for pretrained models while RLHF policies needed a temperature fix spot-checked on three evaluations; neither peer-knowledge prediction nor calibration of a production assistant model is established.
\end{itemize}
